\documentclass[11pt]{article}

\usepackage[a4paper,margin=1in]{geometry}
\usepackage{amsmath,amssymb,amsthm,amscd}
\usepackage{enumitem}
\usepackage[hidelinks]{hyperref}

\title{Address Morphism Theory: A Formal Core}
\author{Address Morphism Theory Project}
\date{Draft v0.1 -- 2026-06-29}

\newtheorem{definition}{Definition}
\newtheorem{assumption}{Assumption}
\newtheorem{theorem}{Theorem}
\newtheorem{proposition}{Proposition}
\newtheorem{corollary}{Corollary}

\newcommand{\Surf}{\mathsf{S}}
\newcommand{\Tok}{\mathsf{T}}
\newcommand{\Ent}{\mathsf{E}}
\newcommand{\Cand}{\mathsf{C}}
\newcommand{\Out}{\mathsf{O}}
\newcommand{\Hist}{\mathsf{H}}
\newcommand{\Ctx}{\mathsf{X}}
\newcommand{\PID}{\mathsf{PID}}
\newcommand{\AGID}{\mathsf{AGID}}
\newcommand{\AOID}{\mathsf{AOID}}
\newcommand{\Resolved}{\mathsf{resolved}}
\newcommand{\Ambiguous}{\mathsf{ambiguous}}
\newcommand{\Unresolved}{\mathsf{unresolved}}
\newcommand{\Rejected}{\mathsf{rejected}}
\newcommand{\Conditional}{\mathsf{conditional}}
\newcommand{\Evidence}{\mathsf{Ev}}
\newcommand{\View}{\mathsf{V}}
\newcommand{\Opt}{\operatorname{Opt}}
\newcommand{\supp}{\operatorname{supp}}
\newcommand{\Pow}{\mathcal{P}}

\begin{document}
\maketitle

\begin{abstract}
Address Morphism Theory (AMT) models an address as a context-indexed,
evidence-bound partial morphism from surface expressions to typed referents.
The goal is not to prove that every address can be resolved. The goal is to
state when resolution is possible, when a system must abstain, and why false
precision is a theoretical error rather than a mere implementation defect.
This short note extracts the formal core of AMT from the larger manuscripts.
\end{abstract}

\section{Novelty Claim}

Most address systems treat an address as one of the following: a formatted
postal string, a coordinate, a postal-code lookup, or a geocoding result. AMT
instead treats address resolution as a chain of typed morphisms under context,
time, evidence, and disclosure policy. Ambiguity, unresolved states, lineage,
and gate-controlled identifier issuance are not edge cases. They are part of
the model.

\section{Objects}

\begin{definition}[Surface expression]
For time \(t\), let \(\Surf_t\) be the set of surface address expressions:
postal strings, partial names, facility names, postal codes, coordinates,
grid labels, \(\AGID\)-like public references, or \(\AOID\)-like private
operation references.
\end{definition}

\begin{definition}[Typed referent]
For context \(\chi \in \Ctx\) and time \(t\), let
\(\Ent_{\chi,t}\) be a typed set of addressable referents. A referent may be a
unit, building, parcel, entrance, road segment, bridge, port, locker, island,
mountain, cultural site, temporary shelter, virtual town, or administrative
object.
\end{definition}

\begin{definition}[Resolution outcome]
Let
\[
\Out_{\chi,t}
  = \Ent_{\chi,t}
    \cup \{\Conditional,\Ambiguous,\Unresolved,\Rejected\}.
\]
The non-final states are first-class outcomes. A resolver that cannot safely
select a referent must return one of them rather than emit a false identifier.
\(\Conditional\) denotes a bounded operational answer such as ``deliverable
only through this entrance'', ``handoff only after human review'', or
``valid only until expiry''. It is not the same as a context-free proof of a
unique address referent.
\end{definition}

\section{Morphism Chain}

AMT models resolution as a partial, evidence-sensitive chain:
\[
\Surf_t
  \xrightarrow{\pi_t}
\Tok_t
  \xrightarrow{\varepsilon_t}
\Pow(\Ent_{\chi,t})
  \xrightarrow{\kappa_{\chi,t}}
\Cand_{\chi,t}
  \xrightarrow{\rho_{\chi,t}}
\Out_{\chi,t}.
\]

\begin{itemize}[leftmargin=1.5em]
\item \(\pi_t\) parses a surface expression into typed tokens.
\item \(\varepsilon_t\) expands aliases, languages, historical names, and
nearby candidates.
\item \(\kappa_{\chi,t}\) assigns evidence, source, lineage, and risk
information.
\item \(\rho_{\chi,t}\) resolves, abstains, or rejects under context \(\chi\).
\end{itemize}

\section{Conditional Commutative Diagrams}

AMT diagrams are conditional. They commute only when the required evidence,
candidate coverage, source freshness, and gate predicates hold.

\subsection{Observation and Partial Inversion}

\[
\begin{CD}
\Ent_{\chi,t} @>{\operatorname{obs}_{\chi,t}}>> \Surf_t \\
@V{\operatorname{id}}VV @VV{\rho_{\chi,t}\circ\kappa_{\chi,t}\circ\varepsilon_t\circ\pi_t}V \\
\Ent_{\chi,t} @<{\operatorname{select}}<< \Out_{\chi,t}
\end{CD}
\]

This square is not assumed to commute for all entities. If
\(\operatorname{obs}_{\chi,t}\) is non-injective, a total perfect inverse cannot
exist.

\subsection{Gate-Controlled Identifier Issuance}

\[
\begin{CD}
\Cand_{\chi,t} @>{G_{\chi,t}}>> \{\mathsf{accept},\mathsf{reject}\} \\
@V{\rho_{\chi,t}}VV @VV{\iota_{\chi,t}}V \\
\Out_{\chi,t} @>{I_{\chi,t}}>> \PID \cup \{\mathsf{no\_emit}\}
\end{CD}
\]

Identifier issuance is permitted only when the certified gate \(G_{\chi,t}\)
accepts the candidate and the resolution outcome is sufficiently precise.

\section{Assumptions}

\begin{assumption}[Finite operational candidates]
For a concrete resolution request, the system evaluates a finite candidate set.
\end{assumption}

\begin{assumption}[Explicit non-emission]
If candidate coverage, source quality, freshness, context, or risk gates are
insufficient, the system may output \(\Ambiguous\), \(\Unresolved\), or
\(\Rejected\).
\end{assumption}

\begin{assumption}[Declared context]
Every resolution claim is relative to a context \(\chi\), such as delivery,
emergency response, property administration, hotel handoff, locker pickup, or
private proof.
\end{assumption}

\section{Core Theorems}

\begin{theorem}[Address reference impossibility]
Suppose there exist distinct referents \(e_1,e_2 \in \Ent_{\chi,t}\) such that
\[
\operatorname{obs}_{\chi,t}(e_1)=\operatorname{obs}_{\chi,t}(e_2).
\]
Then no deterministic resolver
\[
R:\Surf_t\to\Ent_{\chi,t}
\]
can be both total and correct for every referent under this observation map.
\end{theorem}

\begin{proof}
Let \(s=\operatorname{obs}_{\chi,t}(e_1)=\operatorname{obs}_{\chi,t}(e_2)\).
Since \(R\) is a function, \(R(s)\) has one value. Correctness for \(e_1\)
requires \(R(s)=e_1\). Correctness for \(e_2\) requires \(R(s)=e_2\). This
contradicts \(e_1\ne e_2\). Therefore such a resolver cannot exist.
\end{proof}

\begin{corollary}[Unresolved is necessary]
If observation, normalization, candidate generation, or projection can collapse
distinct referents, a safe resolver needs a non-emitting outcome.
\end{corollary}

\begin{proposition}[Lineage is relational]
Administrative splits and merges cannot in general be represented by a single
successor function \(f:\Hist_t\to\Hist_{t+1}\). They require a relation
\[
L_t \subseteq \Hist_t \times \Hist_{t+1}.
\]
\end{proposition}

\begin{proof}
If one historical address state splits into two successor states, a single
function can choose only one successor for that input. The other successor is
lost. A relation can preserve both.
\end{proof}

\section{Context-Relative Optimality}

\begin{definition}[Context utility]
For a surface expression \(s\), context \(\chi\), time \(t\), and candidate
referent \(e \in \Ent_{\chi,t}\), let
\[
U_{\chi,t}(e,s) \in \mathbb{R}
\]
be the operational utility of returning \(e\). The utility may encode delivery
success, emergency access, legal notice, hotel handoff, locker pickup, privacy
risk, or a weighted combination of these terms.
\end{definition}

\begin{definition}[Context-optimal candidate set]
Given a non-empty candidate set \(C_{\chi,t}(s)\subseteq\Ent_{\chi,t}\), define
\[
\Opt_{\chi,t}(s)
  = \arg\max_{e\in C_{\chi,t}(s)} U_{\chi,t}(e,s).
\]
If \(|\Opt_{\chi,t}(s)|>1\), the resolver has a context-level tie and must not
claim unique optimality.
\end{definition}

\begin{theorem}[No context-free operational optimum]
Suppose there are two contexts \(\chi_1,\chi_2\), one surface expression \(s\),
and two candidates \(e_a,e_b\) such that
\[
U_{\chi_1,t}(e_a,s)>U_{\chi_1,t}(e_b,s)
\quad\text{and}\quad
U_{\chi_2,t}(e_b,s)>U_{\chi_2,t}(e_a,s).
\]
Then there is no single context-free resolver that is operationally optimal for
both contexts on \(s\).
\end{theorem}

\begin{proof}
A deterministic context-free resolver must return one candidate for \(s\). If
it returns \(e_a\), it is not optimal under \(\chi_2\). If it returns \(e_b\),
it is not optimal under \(\chi_1\). Therefore operational optimality is a
context-indexed claim.
\end{proof}

\section{Entropy and Evidence}

\begin{definition}[Address reference entropy]
Let \(p_{\chi,t}(e\mid s)\) be the posterior candidate distribution over
\(C_{\chi,t}(s)\). The address reference entropy of \(s\) is
\[
H_{\chi,t}(s)
  =
  -\sum_{e\in C_{\chi,t}(s)}
    p_{\chi,t}(e\mid s)\log p_{\chi,t}(e\mid s).
\]
Entropy measures referential uncertainty, not truth. Low entropy can still be
wrong if the evidence pipeline is biased or incomplete.
\end{definition}

\begin{proposition}[Evidence may reduce uncertainty, not define truth]
Let \(\Evidence\) be an evidence update that changes the posterior from
\(p_{\chi,t}\) to \(p'_{\chi,t}\). If \(\Evidence\) increases the posterior mass
of one correct candidate without adding new collisions or source conflicts,
then \(H'_{\chi,t}(s) < H_{\chi,t}(s)\) may hold. However, entropy reduction
alone does not imply correctness.
\end{proposition}

\begin{proof}
The entropy statement follows from concentrating posterior mass under the
stated condition. Correctness does not follow from entropy alone because a
posterior may concentrate on an incorrect candidate when candidate generation,
source trust, or normalization is wrong.
\end{proof}

\begin{theorem}[Abstention monotonicity]
For fixed \(s,\chi,t\), if the evidence set, candidate set, gate policy, and
freshness metadata are unchanged, then a safe resolver that outputs
\(\Unresolved\), \(\Ambiguous\), or \(\Rejected\) must not later output
\(\Resolved\) merely because time has passed.
\end{theorem}

\begin{proof}
Without new evidence, a changed candidate set, a changed policy, or refreshed
metadata, the proof obligation for safe emission is identical. Emitting a
resolved identifier after abstention would therefore introduce precision not
supported by any new premise.
\end{proof}

\section{Privacy Predicates}

\begin{definition}[Disclosure view]
A disclosure view is a partial map
\[
\View_{\chi,t}:\Ent_{\chi,t}\rightharpoonup D
\]
that reveals only the fields needed for a purpose, such as country eligibility,
delivery zone membership, or locker handoff status.
\end{definition}

\begin{definition}[Predicate anonymity set]
For a predicate \(P:\Ent_{\chi,t}\to\{0,1\}\), define the anonymity set
\[
A_P=\{e\in \Ent_{\chi,t}:P(e)=1\}.
\]
\end{definition}

\begin{proposition}[Private proof granularity]
If \(|A_P|=1\), then proving \(P(e)=1\) identifies the referent within the
modeled universe even if the proof system hides the raw address string. A
privacy-preserving deployment therefore needs a declared minimum anonymity
threshold \(k\) and must treat predicates with \(|A_P|<k\) as unsafe for public
proof.
\end{proposition}

\begin{proof}
When the satisfying set has one element, the predicate value uniquely selects
that element. Cryptography can hide the witness representation, but it cannot
change the information disclosed by the predicate itself.
\end{proof}

\section{Executable Artifact Boundary}

The repository should include a local executable model that checks the minimal
negative cases:

\begin{enumerate}[leftmargin=1.5em]
\item non-injective observation creates ambiguity;
\item missing candidates create unresolved;
\item quality gate failure prevents identifier emission;
\item lineage split is non-functional;
\item no-postcode regions can still be handled as evidence-bound referents.
\item context changes can change the optimal referent;
\item evidence can reduce ambiguity without proving global truth;
\item abstention is monotone without new evidence;
\item privacy predicates require minimum anonymity sets.
\end{enumerate}

These checks do not prove global geographic coverage. They prove that the
repository contains a concrete implementation of AMT's abstention discipline.

\section{Non-Goals}

AMT does not replace postal standards, cadastral law, national address
databases, geocoding APIs, GIS engines, delivery-carrier contracts, or
cryptographic proof systems. It supplies a formal layer that says how these
systems can be composed without pretending that every address is a unique,
stable, globally resolvable string.

\end{document}
